%----------------------------------------------------------------------------------------
%	RESUMEN
%----------------------------------------------------------------------------------------
\cleardoublepage
\thispagestyle{fancy}
\phantomsection
\addtotoc{Resumen}

\begin{center}\huge \textbf{Resumen}\end{center}

\vspace{1em}

\textit{[Resumen de 300 palabras en castellano, siguiendo el formato IMRyD. Espera a terminar el proyecto para elaborarlo. Desde la introducción a las conclusiones describe el proyecto que has realizado. Apartado por apartado, sin nombrarlos.]}

\vspace{2em}

\noindent\textbf{Palabras clave:} \textit{[5 palabras clave para que encuentren tu trabajo cuando forme parte del repositorio de la UAX. Sirven para clasificar tu trabajo en categorías generales.]}

\cleardoublepage

%----------------------------------------------------------------------------------------
%	ABSTRACT (English)
%----------------------------------------------------------------------------------------
\thispagestyle{fancy}
\phantomsection
\addtotoc{Abstract}

\begin{center}\huge \textbf{Abstract}\end{center}

\vspace{1em}

\textit{[Resumen en inglés.]}

\vspace{2em}

\noindent\textbf{Keywords:} \textit{[Palabras clave en inglés.]}

\vspace{2em}

\noindent\small{* El número mínimo de palabras del TFM es de 50\,000 y de 60\,000 como máximo.}

\cleardoublepage
